\chapter{LATAR BELAKANG}

\section{Latar Belakang Masalah}
Kemampuan CNN dalam melakukan ekstraksi fitur hirarkis pada spatial
data matriks menjadikannya sebagai salah satu metode yang paling banyak
digunakan dalam proses pengolahan citra dan visi. Proses ekstraksi fitur
pada data gambar berskala besar dilakukan oleh secara otomatis dan
menjadikan CNN sebagai fondasi utama dalam implementasi pengolahan citra
dan visi untuk pengenalan pola, objek hingga karakter
\cite{Elharrouss2022BackbonesReviewFE,AlexNetPaper}. Keunggulan ini
mendorong adopsi dan implementasi di berbagai bidang terutama sistem tertanam.
Saat ini implementasi CNN banyak dilakukan di GPU dan CPU. Konsumsi daya pada GPU
sangat besar dan boros baik pada saat proses pelatihan model maupun pada saat mode inferensial\cite{GPUBOROS, GPUBATCHINFERENCE}.
Hal ini terjadi terutama pada GPU, karena overhead komunikasi data bus pada batch berukuran kecil sehingga penggunaan daya
sangat tidak efisien.

Keterbatasan implementasi CNN pada CPU dan GPU menjadi semakin
penting ketika CNN diterapkan pada sistem tertanam yang membutuhkan
pemrosesan citra secara langsung, cepat, dan hemat daya\cite{Grigorescu2020AutonomousDrivingSurvey}. Pada aplikasi
seperti kamera cerdas, robotika, dan kendaraan otonom, data visual
diperoleh secara terus-menerus dari sensor sehingga proses komputasi
tidak cukup hanya akurat, tetapi juga harus memiliki latensi rendah.
Oleh karena itu, pendekatan \emph{edge computation} dan
\emph{embedded vision} menjadi salah satu solusi karena proses
inferensi dapat dilakukan lebih dekat dengan sumber data.

Salah satu alternatif perangkat keras yang dapat digunakan untuk
mempercepat komputasi CNN adalah FPGA. FPGA memiliki keunggulan karena arsitektur perangkat
kerasnya dapat dikonfigurasi sesuai kebutuhan algoritma.
Berbeda dengan CPU dan GPU yang menggunakan arsitektur
komputasi umum, FPGA memungkinkan perancang untuk
membangun jalur data khusus, menerapkan paralelisme pada
tingkat operasi, serta menggunakan teknik \emph{pipeline} untuk
meningkatkan \emph{throughput}. Karakteristik ini sesuai dengan
operasi CNN yang bersifat berulang, terstruktur, dan memiliki
pola komputasi yang dapat dipetakan ke dalam perangkat keras pola
\cite{Guo2019FPGANNAcceleratorSurvey,Abdelouahab2018CNNFPGASurvey}.

Penelitian ini berfokus pada perancangan dan implementasi akselerator CNN berbasis
FPGA dengan komputasi \emph{fixed-point} yang berjalan dengan daya
rendah. Akselerator dirancang untuk menjalankan
proses inferensi CNN secara perangkat keras dengan memanfaatkan
paralelisme, \emph{pipeline}, dan optimasi representasi data.
Selain itu, penelitian ini juga membahas integrasi sistem dengan
perangkat kamera dan jalur tampilan citra, sehingga sistem tidak
hanya diuji sebagai modul komputasi terpisah, tetapi juga sebagai
bagian dari sistem \emph{embedded vision} yang menerima data visual
dari lingkungan nyata.

\section{Rumusan Masalah}
Implementasi CNN pada sistem tertanam tidak hanya menuntut akurasi klasifikasi,
tetapi juga efisiensi komputasi, latensi rendah, konsumsi daya yang kecil, serta
kemampuan integrasi dengan perangkat akuisisi dan tampilan citra. CPU dan GPU
memiliki fleksibilitas tinggi, tetapi konsumsi daya dan pola eksekusinya tidak
selalu sesuai untuk sistem \emph{embedded vision} yang bekerja secara langsung
di sisi perangkat. FPGA menjadi alternatif karena jalur data, paralelisme, dan
representasi numeriknya dapat dirancang sesuai kebutuhan inferensi CNN. Namun,
implementasi pada FPGA harus mempertimbangkan keterbatasan DSP, LUT, FF, BRAM,
akses memori, serta sinkronisasi antarperiferal.

Berdasarkan latar belakang tersebut, rumusan masalah dalam penelitian ini adalah
sebagai berikut.
\begin{enumerate}
\item Bagaimana merancang akselerator CNN berbasis komputasi \emph{fixed-point} 16-bit pada FPGA untuk klasifikasi citra MNIST?
\item Bagaimana mengintegrasikan kamera OV7670, DDR SDRAM sebagai \emph{frame buffer}, VGA Pmod Digilent, ROI, CNN accelerator, dan argmax menjadi sistem \emph{embedded vision} end-to-end pada board Arty A7-100T?
\item Bagaimana performa sistem yang dihasilkan dari sisi akurasi, latensi inferensi, penggunaan \emph{resource} FPGA, estimasi konsumsi daya, timing, dan kemampuan kerja \emph{real-time}?
\end{enumerate}

\section{Tujuan Penelitian}
Tujuan penelitian ini adalah merancang, mengimplementasikan, dan mengevaluasi
sistem akselerator CNN berbasis FPGA yang terintegrasi dengan subsistem kamera,
memori, dan tampilan. Secara khusus, tujuan penelitian ini adalah sebagai
berikut.
\begin{enumerate}
\item Merancang akselerator CNN dengan representasi \emph{fixed-point} 16-bit untuk inferensi citra MNIST pada FPGA.
\item Mengintegrasikan subsistem OV7670, DDR SDRAM, VGA Pmod Digilent, ROI, CNN accelerator, dan argmax dalam satu alur sistem \emph{embedded vision}.
\item Mengevaluasi akurasi model sebagai acuan kebenaran fungsi inferensi.
\item Mengevaluasi latensi dan \emph{throughput} inferensi CNN accelerator.
\item Mengevaluasi penggunaan \emph{resource} FPGA, meliputi LUT, FF, BRAM, dan DSP.
\item Mengevaluasi hasil timing dan estimasi konsumsi daya sistem terintegrasi.
\item Mengevaluasi ketercapaian sistem terhadap kebutuhan pemrosesan citra secara \emph{real-time} sesuai hasil implementasi yang diperoleh.
\end{enumerate}

\section{Batasan Masalah}
Batasan masalah pada penelitian ini adalah:
\begin{enumerate}
\item Pada penelitian ini \emph{board} FPGA yang digunakan adalah ARTY 7 100T.
\item Kamera yang digunakan adalah OV7670.
\item Streaming vidio dilakukan lewat VGA, dengan modul VGA PMOD Digilent.
\item Keterbatasan \emph{resource} pada FPGA memungkinkan penelitian menggunakan model yang tidak terlalu besar.
\item \emph{Dataset} yang digunakan adalah \emph{Dataset} MNIST, dengan ukuran 28x28 \emph{pixel} \emph{channel} 1.
\item \emph{Training} model dilakukan di CPU python.
\end{enumerate}


\section{Manfaat Penelitian}
Penelitian ini diharapkan dapat memberikan kontribusi dalam pengembangan implementasi CNN pada perangkat FPGA, khususnya pada penggunaan komputasi \emph{fixed-point}, pemetaan paralelisme operasi konvolusi, dan penerapan \emph{pipeline} untuk mempercepat proses inferensi. Selain itu, penelitian ini dapat menjadi referensi dalam memahami proses perancangan akselerator CNN berbasis perangkat keras untuk aplikasi \emph{embedded vision}.

\section{Sistematika Penulisan}
\noindent
\textbf{BAB I : PENDAHULUAN}

Pada bab ini dijelaskan latar belakang, rumusan masalah, batasan, tujuan, manfaat, dan sistematika penulisan.\\

\noindent
\textbf{BAB II : TINJAUAN PUSTAKA DAN LANDASAN TEORI}

Pada bab ini dijelaskan teori-teori dan penelitian terdahulu yang digunakan sebagai acuan dan dasar dalam penelitian.\\

\noindent
\textbf{BAB III : METODOLOGI PENELITIAN}

Pada bab ini dijelaskan metode yang digunakan dalam penelitian meliputi langkah kerja, pertanyaan penilitian, alat dan bahan, serta tahapan dan alur penelitian.\\

\noindent
\textbf{BAB IV : HASIL DAN PEMBAHASAN}

Pada bab ini dijelaskan hasil penelitian dan pembahasannya.\\

\noindent
\textbf{BAB V : KESIMPULAN DAN SARAN}

Pada bab ini ditulis kesimpulan akhir dari penelitian dan saran untuk pengembangan penelitian selanjutnya.\\